\documentclass[11pt]{article}
\usepackage[T1]{fontenc}
\usepackage{lmodern,amsmath,amssymb,amsthm,mathtools,booktabs,geometry,hyperref,microtype}
\geometry{letterpaper,margin=1.05in}
\hypersetup{colorlinks=true,linkcolor=blue!55!black,urlcolor=blue!55!black,citecolor=blue!55!black}
\usepackage{xcolor}
\newtheorem{theorem}{Theorem}[section]
\newtheorem{lemma}[theorem]{Lemma}
\newtheorem{corollary}[theorem]{Corollary}
\newtheorem{proposition}[theorem]{Proposition}
\theoremstyle{remark}\newtheorem{remark}[theorem]{Remark}
\newcommand{\R}{\mathbb R}\newcommand{\C}{\mathbb C}
\newcommand{\F}{\mathcal F}\newcommand{\dd}{\,\mathrm d}
\newcommand{\ord}{\operatorname{ord}}\newcommand{\Rea}{\operatorname{Re}}
\newcommand{\Ima}{\operatorname{Im}}\newcommand{\sinc}{\operatorname{sinc}}
\title{The sharp endpoint for full-spectrum positive Dirichlet zero filters\\[4pt]\large An endpoint rigidity theorem and exact examples}
\author{A research extension of the ProveIt GRH program}
\date{27 November 2026}
\begin{document}
\maketitle
\begin{abstract}
Assume the Riemann hypothesis for the primitive Dirichlet functions in a finite conjugation-closed multiset $\F$ of size $d$. Put $Q_*=(\prod_{\chi\in\F}q_\chi)^{1/d}$. A positive probability measure cancelling their entire noncentral zero divisor, with multiplicities, cannot have finite moment $\int\exp(b\exp(2|t|/d))\dd\mu(t)$ for $b>\pi d/Q_*$. This halves the earlier necessary bound $2\pi d/Q_*$ in ProveIt and answers its explicit sharpness question in the negative. More strongly, at $b=\pi d/Q_*$ every admissible bilateral Laplace transform is the central-zero-removed product of completed $L$-functions times a polynomial. An exact degree restriction and a positivity test classify all endpoint measures, even asymmetric ones. For every multiset of the zeta function and the primitive quadratic characters of conductors $4$ and $5$, the new bound is the precise supremum. We determine exactly when it is attained, including examples with unique attainment and a transition to nonuniqueness. The proof uses factorization, Stirling's formula, and a sector maximum principle rather than estimates for the spacing of zeros. No new Riemann hypothesis or Artin holomorphy is asserted.
\end{abstract}

\section{Problem and main results}
The starting point is the ProveIt article \emph{Beyond the Riemann hypothesis for Dirichlet L-functions: consequences, transfers and limits}~\cite{proveit}. In its section on full-spectrum positive filters, the result labelled \texttt{fil:thm:global} establishes the optimal double-exponential \emph{exponent} $2/d$. The next corollary, \texttt{fil:cor:endpoint}, gives an upper bound $2\pi d/Q_*$ on the moment coefficient at this exponent, by Jensen's formula. Its open-question section explicitly asks if this coefficient is sharp. The present paper improves that bound and determines the endpoint structure; it does not require new information about individual ordinates or their independence.

Let $\F$ be a nonempty finite multiset of primitive Dirichlet characters, closed under complex conjugation with multiplicity. Include the conductor-one character when desired, denoting its completed function by $\xi$. Write
\begin{align*}
 d&=|\F|, &Q_*&=\Bigl(\prod_{\chi\in\F}q_\chi\Bigr)^{1/d},\\
 \alpha&=2/d, & b_*&=\pi d/Q_*.
\end{align*}
For a nontrivial primitive character of parity $a_\chi\in\{0,1\}$, use
\[
 \Lambda_\chi(s)=(q_\chi/\pi)^{(s+a_\chi)/2}
   \Gamma((s+a_\chi)/2)L(s,\chi),
\]
and for the conductor-one character use $q_\chi=1$, $a_\chi=0$ and
\[
 \Lambda_\chi(s)=\xi(s)=\tfrac12s(s-1)\pi^{-s/2}\Gamma(s/2)\zeta(s).
\]
Define the entire functions and constants
\begin{align}
 F(w)&=\prod_{\chi\in\F}\Lambda_\chi(1/2+w),&
 m&=\ord_{w=0}F,\label{eq:F}\\
 B(w)&=F(w)/w^m,&
 k&=2r-\frac d4+\frac12\sum_{\chi\in\F}a_\chi-m,\label{eq:k}
\end{align}
where $r$ is the number of conductor-one factors and the quotient is extended at zero. Here $k$ is a real exponent, not necessarily an integer. The functional equation and conjugation give $B(-w)=B(w)$ and $B(\bar w)=\overline{B(w)}$. By the assumed critical-line theorem and the usual properties of primitive completed functions, the zero divisor of $B$ is precisely the multiset of nonzero points $i\gamma$ with multiplicities $m(\gamma)=\sum_\chi\ord_{s=1/2+i\gamma}\Lambda_\chi(s)$. In particular $B(x)>0$ for real $x$: it is positive for large positive $x$ by the gamma factors and Euler products, and it has no real zeros.

For a Borel probability measure $\mu$ on $\R$ put
\[
 H(w)=\int_\R e^{wt}\dd\mu(t),\qquad \widehat\mu(y)=H(iy),\qquad
 I_b(\mu)=\int_\R\exp(b e^{\alpha|t|})\dd\mu(t).
\]
The requirement $I_b(\mu)<\infty$ for some $b>0$ makes $H$ entire. Call $\mu$ a \emph{full-divisor filter} if $H$ vanishes at each $i\gamma\ne0$ to at least $m(\gamma)$. This is a stronger requirement than vanishing once at each distinct ordinate; the distinction is essential.

\begin{theorem}[Universal endpoint bound and rigidity]\label{thm:main}
Assume Dirichlet GRH for the characters in $\F$. For any full-divisor filter, $I_b(\mu)<\infty$ implies
\[
 b\le b_*=\frac{\pi d}{Q_*}.
\]
If $I_{b_*}(\mu)<\infty$, then necessarily
\begin{equation}\label{eq:polynomial}
 H(w)=P(w)B(w),\qquad P\in\R[w],\quad P(0)B(0)=1,\quad
 \deg P+k<-\tfrac12.
\end{equation}
Moreover $P(x)>0$ for $x\in\R$, so its degree is even and its leading coefficient is positive. No symmetry of $\mu$ is assumed.
\end{theorem}

\begin{theorem}[Exact characterization of endpoint measures]\label{thm:class}
With the same hypotheses, let
\[
 h(t)=\frac1{2\pi}\int_{\R}B(iy)e^{-iyt}\dd y.
\]
The endpoint measures are precisely the densities
\begin{equation}\label{eq:density}
 p(t)=P(-\partial_t)h(t)\ge0\quad(t\in\R)
\end{equation}
for real polynomials $P$ satisfying $P(0)B(0)=1$ and $\deg P+k<-1/2$. Thus if $k\ge-1/2$ the endpoint is impossible; if $-5/2\le k<-1/2$ there is at most one endpoint measure, namely $h(t)\dd t/B(0)$, and it exists exactly when $h\ge0$.
\end{theorem}

The last assertion follows because a real polynomial positive on the whole real axis has even degree, and a positive degree would have to be at least two. Notice that Theorem~\ref{thm:class} is a finite-dimensional criterion, not a claim that the signed kernel $h$ is always positive. For arbitrary families we do not assert that $b_*$ is a sharp supremum; the previous existence theorem guarantees some positive $b$. The examples in Section~\ref{sec:examples} prove sharpness on a large concrete class.

\section{Growth and the entire quotient}
\begin{lemma}[Sector asymptotics]\label{lem:stirling}
Uniformly for $w=R e^{i\theta}$, $|\theta|\le\pi/4$, as $R\to\infty$, with $x=\Rea w$ and $y=\Ima w$,
\begin{equation}\label{eq:stirling}
 \log|B(w)|=\frac d2\left(x\log R-y\theta-x+x\log\frac{Q_*}{2\pi}\right)+k\log R+O_\F(1).
\end{equation}
In particular there is a constant $C_\F>0$ with
\begin{equation}\label{eq:realasymp}
 B(x)\sim C_\F x^k\exp\left[\frac x\alpha\left(\log\frac{x}{b_*\alpha}-1\right)\right]
 \quad(x\to+\infty).
\end{equation}
\end{lemma}
\begin{proof}
In this sector $\Rea(1/2+w)\to+\infty$, so every Dirichlet series is $1+o(1)$ uniformly by absolute convergence (the leading coefficient is $1$). With $c_\chi=1/4+a_\chi/2$, sectorial Stirling gives
\[
 \Gamma(w/2+c_\chi)=\sqrt{2\pi}(w/2)^{w/2+c_\chi-1/2}e^{-w/2}(1+O_\F(R^{-1})).
\]
The log modulus of the main power is computed with the principal logarithm. Multiplying by $(q_\chi/\pi)^{w/2+c_\chi}$ contributes $\tfrac12x\log(q_\chi/\pi)+O_\F(1)$. Each zeta factor also has the polynomial $s(s-1)/2$, contributing $2\log R+O(1)$; division by $w^m$ subtracts $m\log R$. Summation proves \eqref{eq:stirling}, and retaining the leading positive constants on the real axis proves \eqref{eq:realasymp}, since $b_*\alpha=2\pi/Q_*$.
\end{proof}

\begin{lemma}[Divisibility and growth]\label{lem:division}
If $I_b(\mu)<\infty$ and $\mu$ is a full-divisor filter, then $G=H/B$ is entire and for every $\sigma>1$ it satisfies $\log^+|G(w)|=O_{\sigma,\mu,\F}(1+|w|^\sigma)$.
\end{lemma}
\begin{proof}
For $u=|\Rea w|\ge b\alpha$ one has, by maximizing $uv-be^{\alpha v}$ for $v\ge0$,
\begin{equation}\label{eq:legendre}
 |H(w)|\le I_b(\mu)\exp\left[\frac u\alpha\left(\log\frac{u}{b\alpha}-1\right)\right].
\end{equation}
For smaller $u$ the maximum is attained at $v=0$, giving a bounded estimate. These estimates also justify differentiation under the integral locally uniformly and give $\log^+ M_H(R)=O(R\log(2+R))$. The completed functions have order at most one (classical; see~\cite{davenport,titchmarsh}); for example their theta-integral representations have kernels bounded by $C\exp(-c e^{2|t|})$, by the theta transformation, and give the same growth bound. Finite products and removal of a central zero preserve order at most one. GRH identifies the entire zero divisor of $B$ with the required noncentral divisor. Thus $H/B$ has removable singularities everywhere.

It remains to bound its order, which cannot be inferred from a naive pointwise lower bound near the zeros. Hadamard factorization~\cite{boas} applies to $H$ and $B$, both nonzero at the origin, with canonical genus-one factors $E_1(z)=(1-z)e^z$ and exponential prefactors of degree at most one. Cancel the divisor of $B$ in the absolutely locally convergent canonical products to obtain
\[
 G(w)=e^{Aw+C}\prod_{z\in Z_{\rm extra}} E_1(w/z),
\]
where $Z_{\rm extra}$ is the surplus multiset of zeros of $H$. For $1<\sigma<2$, $\sum_{z\in Z_{\rm extra}}|z|^{-\sigma}<\infty$ and $\log |E_1(v)|\le C_\sigma|v|^\sigma$ (use the Taylor expansion for $|v|\le1/2$ and the elementary bound for $|v|>1/2$). This proves the asserted growth for $1<\sigma<2$, hence also for $\sigma\ge2$.
\end{proof}

\section{The endpoint quotient is a polynomial}\label{sec:rigidity}
Suppose $I_{b_*}(\mu)<\infty$. Combining \eqref{eq:legendre} with \eqref{eq:stirling} cancels every term of order $R\log R$ and every conductor-dependent linear term. Uniformly in the right cone, for large $R$,
\begin{equation}\label{eq:angular}
 |G(R e^{i\theta})|\le C R^{-k}
 \exp\left[\frac d2 R\{\cos\theta\log\cos\theta+\theta\sin\theta\}\right].
\end{equation}
The same bound applies in the left cone after replacing $w$ by $-w$ in the estimate for $B$ and using $|\Rea w|$ for $H$; this does not assume that $H$ is even. The term $\cos\theta\log\cos\theta$ is nonpositive. Hence for $|\theta|\le\theta_0\le\pi/4$, the exponential in \eqref{eq:angular} is at most $\exp((d/2)\theta_0|y|)$.

\begin{lemma}[Two narrow cones force a polynomial]\label{lem:PL}
Let $G$ be entire with $\log^+|G(w)|=O_\sigma(1+|w|^\sigma)$ for all $\sigma>1$. Suppose that for each $\eta>0$ there are a half-angle $0<\theta_0<\pi/4$ and a constant $C_\eta$ such that in the two cones about the positive and negative real axes of half-angle $\theta_0$,
\[
 |G(w)|\le C_\eta(1+|w|)^N e^{\eta|\Ima w|},
\]
where $N$ is a fixed nonnegative integer. Then $G$ is a polynomial of degree at most $N$.
\end{lemma}
\begin{proof}
Fix $\eta$ and work in the upper half-plane. The holomorphic function
\[
 g_\eta(w)=\frac{G(w)}{(w+i)^N}e^{i\eta w}
\]
is bounded in the two narrow cones, because $|w+i|\ge\max(1,|w|)$ for $\Ima w\ge0$. On the intervening sector, of opening $\pi-2\theta_0<\pi$, it is bounded on the rays and has growth $\exp(O(|w|^\sigma))$ for any $1<\sigma<\pi/(\pi-2\theta_0)$. The sector Phragm\'en--Lindel\"of principle~\cite{boas} therefore bounds $g_\eta$ throughout that sector. Together these bounds give boundedness on the upper half-plane. On the real boundary it is bounded by the constant
\[
 C_0=\sup_{x\in\R}|G(x)/(x+i)^N|<\infty,
\]
independent of $\eta$. For completeness, multiply $g_\eta(w)$ by $(1-i\delta w)^{-1}$ for $\delta>0$; on an upper semicircle of radius $R$, the denominator has modulus at least $\delta R$, and on the real line it is at least one. The maximum principle on half-disks followed by $R\to\infty$, then $\delta\downarrow0$, gives $|g_\eta(w)|\le C_0$ in the entire upper half-plane. Therefore $|G(w)|\le C_0|w+i|^N e^{\eta\Ima w}$. Let $\eta\downarrow0$. The lower half-plane follows with $(w-i)^N$ and $e^{-i\eta w}$. The resulting polynomial growth on the whole plane implies the conclusion by Cauchy's coefficient estimates.
\end{proof}

Apply Lemma~\ref{lem:PL} with $N=\max(0,\lceil-k\rceil)$, choosing $\theta_0\le 2\eta/d$ in \eqref{eq:angular}. This proves $G=P$ is a polynomial. Its real coefficients follow from $H(\bar w)=\overline{H(w)}$ and the same identity for $B$. At zero it satisfies $P(0)B(0)=1$. On the real line $H(x)>0$ because $\mu$ is a probability, and $B(x)>0$ by the critical-line input, so $P(x)>0$. Its leading coefficient is positive and its degree is even.

If $I_b(\mu)<\infty$ with $b>b_*$, then $I_{b_*}(\mu)<\infty$ as well, so $G=P$ is a polynomial. But \eqref{eq:legendre} on the positive real axis together with \eqref{eq:realasymp} now gives
\[
 |P(x)|\ll x^{-k}\exp\left[-\frac x\alpha\log(b/b_*)\right],
\]
which tends to zero exponentially. This contradicts $P(0)B(0)=1$. We have proved the improved upper bound. The argument has not used a numerical zero count or a zero-spacing estimate.

\section{The exact moment criterion at the endpoint}
A pointwise bound on the real Laplace transform alone would miss an important half-power. We next identify it using a positive integral, so the condition in Theorem~\ref{thm:main} is strict and exact.

\begin{lemma}[Saddle integral equivalence]\label{lem:saddle}
Let $a,b>0$, and let $\mu$ be a finite positive measure with all real exponential moments, with bilateral Laplace transform $H$. Define
\[
 f(x)=\frac{x}{a}\left(\log\frac{x}{ab}-1\right)\quad(x>0).
\]
Then
\begin{equation}\label{eq:momenttest}
 \int e^{b e^{a|t|}}\dd\mu(t)<\infty
 \quad\Longleftrightarrow\quad
 \int_1^\infty (H(x)+H(-x))e^{-f(x)}x^{-1/2}\dd x<\infty.
\end{equation}
\end{lemma}
\begin{proof}
Put $J(t)=\int_1^\infty e^{xt-f(x)}x^{-1/2}\dd x$. By Tonelli the right-hand integral is $\int(J(t)+J(-t))\dd\mu(t)$. It suffices to prove
\[
 J(t)+J(-t)\asymp_{a,b}e^{b e^{a|t|}}.
\]
For $t\to+\infty$ put $x_0=ab e^{at}$ and $D(u)=u\log u-u+1$. Substitution $x=x_0u$ gives
\begin{equation}\label{eq:J}
 J(t)e^{-b e^{at}}=\sqrt{x_0}\int_{1/x_0}^\infty u^{-1/2}\exp[-(x_0/a)D(u)]\dd u.
\end{equation}
On $[1/2,2]$, $D(1)=D'(1)=0$ and $1/2\le D''(u)\le2$, hence $(u-1)^2/4\le D(u)\le(u-1)^2$. The Gaussian comparison yields a bounded upper contribution and a positive lower contribution, the latter by integrating on $|u-1|\le x_0^{-1/2}$ for large $x_0$. Below $1/2$, $D(u)\ge D(1/2)>0$ and $u^{-1/2}$ is integrable. Above $2$, $D(u)\ge D(2)+(u-2)\log2$. Both remaining contributions are bounded after multiplication by $\sqrt{x_0}$. Thus \eqref{eq:J} is bounded above and below by positive constants for large $t$. Since $J(-t)\le J(0)<\infty$ for $t\ge0$ and both functions are continuous and positive on bounded intervals, the result follows by reflection and Tonelli.
\end{proof}

Use Lemma~\ref{lem:saddle} with $a=\alpha$, $b=b_*$. If $H=PB$ is the transform of a positive probability, $P$ has even degree $j$ and positive leading coefficient. By \eqref{eq:realasymp} and evenness of $B$,
\[
 (H(x)+H(-x))e^{-f(x)}x^{-1/2}\asymp x^{j+k-1/2}\quad(x\to+\infty).
\]
The integral converges precisely when $j+k-1/2<-1$, namely $j+k<-1/2$. This proves the remaining assertion of Theorem~\ref{thm:main}. It also proves the moment equivalence for \emph{any} positive probability having the transform $PB$, provided all exponential moments exist. At equality of the power to $-1$ the integral diverges logarithmically.

\begin{proof}[Proof of Theorem~\ref{thm:class}]
On every fixed vertical strip, Stirling's formula for the gamma factors and the classical polynomial growth of Dirichlet $L$-functions in fixed strips imply that $B(c+iy)$ is bounded by a polynomial in $1+|y|$ times $\exp(-\pi d|y|/4)$, uniformly for $c$ in bounded intervals. (The finite central division has only a removable singularity.) Hence the inverse Fourier integral defining $h$ and all of its derivatives converge absolutely; $h$ is real and even. For any polynomial $P$, inversion gives the real continuous function $p_P=P(-\partial_t)h$ when $P$ has real coefficients. Shift the contour of the inverse Fourier integral for $PB$ to $\Rea w=c$ for $t>0$, and to $\Rea w=-c$ for $t<0$. The horizontal sides vanish by the uniform vertical-strip estimate. For each fixed $c>0$ this proves $|p_P(t)|\ll_{P,c,\F}e^{-c|t|}$. In particular all real exponential moments of $|p_P|$ exist.

Fourier inversion and uniqueness of Fourier transforms of finite measures now show that if $H=PB$ for a probability measure, that measure has the density $p_P$; nonnegativity holds everywhere by continuity. Conversely, if $p_P\ge0$ and $P(0)B(0)=1$, it is a probability density, its Fourier transform is $P(iy)B(iy)$, and analytic continuation gives its bilateral Laplace transform $PB$. Its required divisor is cancelled and Lemma~\ref{lem:saddle} gives the exact criterion $\deg P+k<-1/2$. This proves both directions.
\end{proof}

\section{Sharpness and attainment for concrete families}\label{sec:examples}
We now establish exact results for every multiset made of the conductor-one character, $\chi_{-4}$ (Dirichlet beta), and $\chi_5$ (the real primitive quadratic character modulo $5$). This also supplies a positive theta kernel not used in the cited global-filter examples.

\begin{lemma}[Three positive kernels]\label{lem:kernels}
Each of the three completions is the bilateral Laplace transform of an even, strictly positive smooth integrable kernel $K_\chi(t)$. For constants $C_\chi,A_\chi>0$,
\[
 K_\chi(t)\le C_\chi e^{A_\chi|t|}\exp[-(\pi/q_\chi)e^{2|t|}].
\]
In particular none of these completed functions has a central zero.
\end{lemma}
\begin{proof}
For $t\ge0$ the classical zeta kernel is
\[
 K_1(t)=\sum_{n\ge1}(4\pi^2n^4e^{9t/2}-6\pi n^2e^{5t/2})e^{-\pi n^2e^{2t}}.
\]
It is even by the theta transformation and gives $\xi(1/2+w)$ by Mellin integration and integration by parts; see~\cite[\S2.16]{titchmarsh} and~\cite[\texttt{fil:lem:kernels}]{proveit}. Every displayed term is positive because $2\pi n^2 e^{2t}>3$. For beta the kernel is
\[
 K_{-4}(t)=2e^{3t/2}\sum_{j\ge0}(-1)^j(2j+1)e^{-\pi(2j+1)^2 e^{2t}/4}.
\]
The terms in absolute value strictly decrease for $t\ge0$ (the derivative of $x e^{-a x^2}$ is negative on $x\ge1$ for $a\ge\pi/4>1/2$), so the alternating sum is positive. The theta transformation with root number $1$ gives evenness.

For the third kernel put $a=\pi e^{2t}/5$ and
\[
 K_5(t)=2e^{t/2}\sum_{n\ge1}\chi_5(n)e^{-\pi n^2e^{2t}/5}.
\]
For $t\ge0$, $a\ge\pi/5>1/2$. The function $g(x)=e^{-a x^2}$ is strictly convex on $x\ge1$, since $g''(x)=2a(2ax^2-1)e^{-ax^2}>0$. The character has values $+1,-1,-1,+1,0$ on successive residues $1,2,3,4,0$. Each block contributes
\[
 g(5j+1)+g(5j+4)-g(5j+2)-g(5j+3)>0,
\]
by strict convexity (the two outer points have the same sum as the inner points). Absolute convergence therefore proves positivity for $t\ge0$. Poisson summation for the primitive even quadratic character gives $\theta_5(x)=x^{-1/2}\theta_5(1/x)$ with root number $1$ (Gauss sum $\sqrt5$), so $K_5(-t)=K_5(t)$. For either nontrivial character termwise Mellin integration on the initial half-plane gives exactly $\Lambda_\chi(1/2+w)$, and theta inversion extends the identity to all $w$.

In each case factor $\exp[-(\pi/q_\chi)e^{2t}]$ out of the series on $t\ge0$ and bound the remaining sum by a convergent Gaussian sum with polynomial coefficients, using $e^{2t}\ge1$. This gives the claimed envelope; evenness handles the other tail. The central values are the positive integrals of these kernels.
\end{proof}

\begin{theorem}[Exact supremum and endpoint attainment]\label{thm:examples}
Suppose $\F$ consists of $r$ copies of the zeta factor, $s$ copies of the beta factor, and $u$ copies of the conductor-$5$ factor, with $d=r+s+u>0$. Under the critical-line input, the supremum of all $b$ for which a full-divisor filter has finite $I_b$ is
\[
 b_*=\frac{\pi(r+s+u)}{(4^s5^u)^{1/d}}.
\]
The supremum is attained if and only if
\begin{equation}\label{eq:attain}
 u>7r+s+2.
\end{equation}
When $7r+s+2<u\le7r+s+10$, there is exactly one endpoint measure. When $u>7r+s+10$, there are infinitely many endpoint measures. All assertions about nonattainment and uniqueness concern arbitrary Borel probability measures, not only even or smooth filters.
\end{theorem}
\begin{proof}
Normalize the positive kernels from Lemma~\ref{lem:kernels} and convolve the $d$ factors. Equivalently let $X_j$ be independent random variables with these normalized densities and set $T=\sum_j X_j$. Its transform is $F(w)/F(0)=B(w)/B(0)$ since $m=0$. For $0<b<b_*$, choose a common $0<\varepsilon<1$ so that $b<(1-\varepsilon)b_*$, and set $c_j=(1-\varepsilon)\pi/q_j$. The kernel envelope gives $\mathbb E\exp(c_j e^{2|X_j|})<\infty$. By the arithmetic--geometric mean inequality and the triangle inequality,
\[
 \sum_{j=1}^d c_j e^{2|X_j|}\ge d(\prod c_j)^{1/d}e^{(2/d)\sum |X_j|}
 \ge (1-\varepsilon)b_*e^{\alpha|T|}.
\]
Independence gives $I_b(\mu_T)\le\prod_j\mathbb E\exp(c_j e^{2|X_j|})<\infty$. Thus every subendpoint coefficient is possible, and Theorem~\ref{thm:main} gives the exact supremum. The polynomial power in \eqref{eq:k} equals $k=(7r+s-u)/4$. The canonical convolution corresponds to the constant polynomial $1/B(0)$, so Theorem~\ref{thm:class} shows that it attains the endpoint precisely if $k<-1/2$. If $k\ge-1/2$, no polynomial of nonnegative degree can qualify. This proves \eqref{eq:attain}. For $-5/2\le k<-1/2$, the degree condition and even degree force a constant polynomial, hence uniqueness.

Finally suppose $k<-5/2$, i.e. $u>7r+s+10$. In particular there is a conductor-$5$ factor. Its positive kernel satisfies $K_5''(t)/K_5(t)\to+\infty$ as $|t|\to\infty$: as $t\to+\infty$, with $X=e^{2t}$, the $n=1$ summand is $2e^{t/2}e^{-(\pi/5)X}$, whose logarithmic derivative is $1/2-(2\pi/5)X$ and whose second derivative divided by itself is $(1/2-(2\pi/5)X)^2-(4\pi/5)X$. The differentiated $n\ge2$ terms are relatively exponentially small even after multiplication by any fixed power of $X$ (factor the first Gaussian and bound the remaining series). Evenness gives the negative tail. Continuity and positivity now give $K_5''\ge-MK_5$ for some finite $M>0$. For every $0<c<1/M$, $K_5+cK_5''>0$. Its integral is unchanged, since the first derivative tends to zero at both ends. Replace one normalized $K_5$ in the convolution by the corresponding normalized $K_5+cK_5''$. The transform is now $(1+cw^2)B(w)/B(0)$, with degree two and $2+k<-1/2$. The endpoint criterion proves finiteness, and different $c$ give different measures by uniqueness of Laplace transforms. This proves nonuniqueness.
\end{proof}

\begin{center}
\small
\begin{tabular}{@{}lrrll@{}}
\toprule
Family & $d$ & $k$ & $b_*$ & Endpoint\\
\midrule
$\zeta$ &1&$7/4$&$\pi$& not attained\\
$L(s,\chi_{-4})$&1&$1/4$&$\pi/4$& not attained\\
$L(s,\chi_5)$&1&$-1/4$&$\pi/5$& not attained\\
$L(s,\chi_5)^2$&2&$-1/2$&$2\pi/5$& not attained\\
$L(s,\chi_5)^3$&3&$-3/4$&$3\pi/5$& unique\\
$L(s,\chi_5)^{10}$&10&$-5/2$&$2\pi$& unique\\
$L(s,\chi_5)^{11}$&11&$-11/4$&$11\pi/5$& infinitely many\\
\bottomrule
\end{tabular}
\end{center}

\section{Consequences, limits, and further questions}
For a smooth prime statistic with an absolutely convergent critical-line expansion $E(t)=\sum_\gamma a_\gamma e^{i\gamma t}$ (or with a separate controlled residual), averaging $E(t+v)$ against a filter gives precisely the central term: by absolute convergence every noncentral term is multiplied by $H(i\gamma)=0$. The repository derives this setting from the smoothed explicit formula. Our theorem constrains the attainable tail of every positive averaging law satisfying the stronger multiplicity-preserving requirement. The exact threshold $b_*$ means that making a filter more concentrated by demanding a larger double-exponential moment coefficient is impossible. This does not constrain signed filters in the same way; positivity was used in the moment estimates, in the saddle integral, and to rule out real zeros of $H$.

The proof uses the repository's transfer strategy in a new way: an upper estimate for a positive analytic object is compared against a lower, asymptotically exact estimate for an arithmetic object with the same required zero divisor. Divisibility and order control make their quotient analytic, the main growth cancels, and a maximum principle supplies rigidity. This is analogous in structure, not identical in implementation, to the source proof's comparison of two bounds followed by a one-sided analytic argument. No Selmer-group or automorphic lifting assertion is needed here.

\paragraph{What is settled.} The explicit sharpness question in \texttt{fil:sec:open} about $2\pi d/Q_*$ is answered negatively in every family: the factor two is unnecessary. The replacement is universally necessary, exactly sharp for the class in Theorem~\ref{thm:examples}, and its endpoint solutions for any family have the necessary-and-sufficient description in Theorem~\ref{thm:class}. The conductor-$5$ positivity proof expands the known example class and demonstrates actual attainment, not only a supremum.

\paragraph{What remains.} For general conjugation-closed families we have not proved that $b_*$ is the supremum: positivity of the inverse kernel or of a suitable multiplier must be established. For $k<-1/2$ the classification reduces endpoint attainment to a finite-dimensional positivity question, but it does not decide it for every family. Requiring cancellation at distinct ordinates only, rather than at full multiplicities, may change the bound. Finite-support filters, including their second-order support asymptotics, are a different optimization problem. The present results do not settle the full nonabelian GRH, Artin's conjectures, or any other major unresolved classical conjecture.

\paragraph{Further research questions.}
\begin{enumerate}
\item For every conjugation-closed Dirichlet family, is the supremum $\pi d/Q_*$ even when the central-zero-removed inverse kernel changes sign? Can one produce positive polynomial modifications or multipliers of arbitrarily small exponential type?
\item Classify all endpoint polynomials in Theorem~\ref{thm:class} for the quadratic families, including asymmetric filters. For high powers of the conductor-$5$ factor, determine the boundary of the cone of positive differential modifications.
\item Transfer the quotient argument to completed automorphic or Dedekind zeta functions with several gamma factors, conditional on their own critical-line statement. The gamma-weighted conductor and the central vanishing order should play the same roles, but their precise normalization must be checked.
\item Study the analogous extremum with one zero per distinct ordinate. Without control of the lost multiplicity divisor, the present division argument cannot be reused verbatim.
\end{enumerate}

\section*{Source record and proof status}
\begingroup\sloppy
The GitHub source was inspected through the repository connector. The inspected \texttt{fil.tex} blob is \path{8f4e47a5e3e1ae4e2381aa3305c938e3906b35c4}, also verified at ProveIt commit \path{a4021de208d592a04419593436847c8e9db9b456}; the upstream input is pinned by the repository at \path{72fead102f663d5c7a407964a13dc0cf9d977207}. See \texttt{SOURCE\_AUDIT.md} in the archive. The repository states the input theorem for Dirichlet functions and describes its article as a merge, not an independent review. We take this precise critical-line theorem as an assumption of the new conclusions; we have not rebuilt the upstream Lean development or independently certified RH/GRH. The mathematical argument above is analytic and requires no computational certificate. The accompanying standard-library Python script verifies the rational exponent and degree bookkeeping for the example classification, not the analytic proof. No new machine-checked Lean proof or global publication-priority claim is made. A targeted literature search found classical kernel and factorization results but no exact match to this endpoint formulation; this is evidence of novelty relative to the inspected sources only.

\endgroup
\begin{thebibliography}{9}
\bibitem{proveit} V. Reshetnikov, ProveIt, \emph{Beyond the Riemann hypothesis for Dirichlet L-functions: consequences, transfers and limits}, repository snapshot a4021de2, Section ``Positive filters that cancel zeros'', especially labels \texttt{fil:lem:kernels}, \texttt{fil:thm:global}, \texttt{fil:cor:endpoint}, \texttt{fil:sec:open}. \href{https://github.com/VladimirReshetnikov/ProveIt/blob/a4021de208d592a04419593436847c8e9db9b456/Analysis/GeneralizedRiemannHypothesis/Article/sections/fil.tex}{Source file}.
\bibitem{upstream} L. Sernau et al., \emph{grh-lean}, commit \path{72fead102f663d5c7a407964a13dc0cf9d977207}, \href{https://github.com/lsernau/grh-lean/tree/72fead102f663d5c7a407964a13dc0cf9d977207}{upstream repository}; used as input by ProveIt.
\bibitem{boas} R. P. Boas, Jr., \emph{Entire Functions}, Academic Press, 1954. (Hadamard factorization and the Phragm\'en--Lindel\"of principle.)
\bibitem{davenport} H. Davenport, \emph{Multiplicative Number Theory}, 3rd ed., revised by H. L. Montgomery, Springer, 2000. (Primitive character theta transformation, functional equation and analytic growth.)
\bibitem{titchmarsh} E. C. Titchmarsh, \emph{The Theory of the Riemann Zeta-Function}, 2nd ed., revised by D. R. Heath-Brown, Oxford, 1986.
\end{thebibliography}
\end{document}
